\chapter{RESULTADOS ESPERADOS}

\orientacao{No TCC I ainda não há resultados empíricos — esta seção descreve o que você ESPERA encontrar, com base na literatura e na metodologia proposta. Seja realista e fundamente cada expectativa.}

\placeholder{Parágrafo introdutório: ``Por se tratar de um Trabalho de Conclusão de Curso I, esta pesquisa encontra-se em fase de planejamento. Nesta seção, apresentam-se os resultados esperados com base na revisão bibliográfica e na metodologia delineada.''}

\placeholder{Resultado esperado 1: descreva o achado principal que você antecipa. Fundamente com citações da literatura (\texttt{\textbackslash cite\{\}}).}

\placeholder{Resultado esperado 2: descreva um segundo achado ou cenário alternativo possível. Mostre que você considerou diferentes desfechos.}

\placeholder{Contribuição acadêmica esperada: como os resultados podem avançar o conhecimento na área?}

\placeholder{Contribuição prática esperada: como os resultados podem ser úteis para profissionais, instituições ou sociedade?}

\orientacao{No TCC II, este capítulo será substituído por ``Resultados'' com dados, tabelas e gráficos reais. Mantenha a mesma estrutura lógica ao migrar o conteúdo.}
